\chapter{Experimental Evidence and Result Artifacts}
\label{app:evidence}

This appendix maps every quantitative result reported in Chapter~\ref{chap:experiments} to the artifact it was read from, so that each number can be verified at its source. The recognition numbers were produced by the build notebook \path{notebooks/Model_D1_Final_Build_Runtime_Contract.ipynb}, which writes its results to \path{notebooks/model_HAR_27may_artifacts/}; this is the same location referenced by the deployed contract metadata.

A note on protocol provenance. The five-fold result uses a stratified window-level split executed in this build notebook. The within-distribution result in this thesis comes from the build notebook, while the pipeline orchestrates inference, monitoring, evaluation, and governed retraining over the resulting runtime contract. The from-scratch architecture in the pipeline training module is a smaller baseline configuration; the deployed proposed model (967{,}499 parameters) was built and fine-tuned in the notebook, and governed retraining fine-tunes that deployed model rather than training a new architecture from scratch.

\renewcommand{\arraystretch}{1.25}
\begin{table}[htbp]
\centering
\footnotesize
\caption[Chapter~7 results and their source artifacts.]{Chapter~7 results and their source artifacts.}
\label{tab:app_evidence}
\begin{tabularx}{\textwidth}{p{3.4cm} p{2.0cm} X}
\toprule
\textbf{Result} & \textbf{Value} & \textbf{Source artifact} \\
\midrule
Five-fold window-level CV & 93.61\% / 0.937 & \path{notebooks/model_HAR_27may_artifacts/garmin_5fold_window_cv_results.json} \\
Blind User~7 probe & 36.96\% / 0.350 / conf 0.858 & \path{models/runtime_contract/model_d1_final_20260528/user7_blind_probe_summary.json} \\
Chapter~7 figures (7.1--7.4) & --- & \path{scripts/generate_ch7_figures.py} \\
Confusion matrix and reliability & --- & regenerated from the deployed model on the User~7 session by the script above \\
Monitoring and trigger outputs (Section~\ref{sec:ch7_monitoring}) & \texttt{review} & \path{thesis/evidence/TRIGGER_POLICY_EVAL.csv}; per-run \path{trigger_decision.json} \\
Contract-sensitivity sweeps (Sections~\ref{sec:ch7_ablation}, \ref{sec:ch7_triggerpolicy}) & --- & \path{thesis/evidence/ABLATION_WINDOWING.csv}; \path{reports/THRESHOLD_CALIBRATION.csv} \\
Unlabeled monitoring at scale (25 sessions, Figures~\ref{fig:ch7_mon_confidence}--\ref{fig:ch7_mon_dominant}) & 23/25 review, 0/25 retrain & \path{scripts/run_unlabeled_sessions.py} and \path{scripts/collect_unlabeled_monitoring.py} produce \path{reports/unlabeled_sessions_monitoring.json} and \path{.csv} from the per-run artifact trails \\
User-7 labeled vs unlabeled (Figure~\ref{fig:ch7_proxy_comparison}) & 36.96\% vs 36.70\% & \path{models/runtime_contract/model_d1_final_20260528/user7_blind_probe_summary.json} and \path{user7_unlabeled_probe_summary.json} \\
Temporal smoothing effect (Section~\ref{sec:ch7_sessions25}, Figure~\ref{fig:ch7_temporal_smoothing}) & $T$: 40.6\%$\to$21.5\%; acc.\ +2.7pp & \path{scripts/generate_smoothing_figure.py} on the User~7 unlabeled session \\
\bottomrule
\end{tabularx}
\renewcommand{\arraystretch}{1.0}
\end{table}

The five-fold result is recorded with the explicit note that it measures within-distribution learnability. These provenance notes are the basis for the two-protocol framing used throughout Chapter~\ref{chap:experiments}.
